% ===============================================================
% FILE: main_paper.tex
% ARTÍCULO PARA REVISTA INDEXADA IEEE (IEEEtran Class)
% Objetivo: Vender el avance en menos de 10 páginas, enfocándose solo en los resultados superiores.
% ===============================================================

\documentclass[conference]{IEEEtran} % CLASE ESTÁNDAR DE IEEE
\usepackage{amsmath, amssymb}
\usepackage{graphicx} 

\title{Graph Signal Processing for Enhanced EEG Data Imputation and Artifact Removal.}
\author{\IEEEauthorblockN{Tu Nombre}
\IEEEauthorblockA{Your Department, Your University \\ \textit{your.email@uni.edu}}}

\begin{document}

\maketitle 

% --- ABSTRACT (Inglés) --- 
\begin{abstract}
(Aquí va el contenido del Abstract que generamos anteriormente: debe ser conciso y contener los términos clave como "spatio-temporal", "Graph Signal Processing," y la superioridad en $\text{MAE}$).
\end{abstract}

\section{Introduction}
(Máximo 1 página). Ir directo al punto. El párrafo de contexto debe terminar con una pregunta que solo tu método puede responder.

\section{Methodology: The GSP-Enhanced Pipeline}
(Aquí se describe el proceso, pero *sin* desviaciones. Se centra en $\mathbf{W}$, las tres familias de tareas, y el protocolo de validación para mostrar la complejidad manejada.) 

\subsection{Graph Construction and Connectivity Weighting}
(Solo se debe incluir la matemática del Kernel Gaussiano vs Data-Driven sin detenerse demasiado en los detalles básicos).

% ... (y así sucesivamente con las secciones más técnicas) ...

\section{Results: Performance Benchmarking and Significance}
Aquí es donde va el peso de tu trabajo. Se usa la tabla $\text{b3\_rep01_final\_table}$. El texto debe guiar al lector a través de los *hallazgos*.

\subsection{Quantitative Superiority of TV/Time Methods}
La Tabla \ref{tab:main_comparison_it82} muestra que el rendimiento es condicional por dominio, destacando la superioridad de $\text{TV/Tiempo}$ en contextos dinámicos. (Mostrar la gráfica clave).

\section{Conclusion}
(Breve y potente). Reiterar los 3 puntos más importantes.

\bibliographystyle{IEEEtran}
\bibliography{tesis_biblio} 

\end{document}
